\section{Introduction}
\label{sec:intro}

Some moves matter for what they make possible next. A Monopoly player in debt who sells the houses on a color group receives the sale price and also gains something the menu does not show, the right to mortgage that group's deeds. A player who sells the last deed in a rival's color group receives cash and gives the rival the right to build. In both cases the terms on the table understate the move. An agent that reads only those terms will quit a debt it could have paid, or sell a monopoly at the price of a single deed.

Language-model agents now play games of this kind with free-form negotiation, and they write down reasons as they act. That makes the problem observable. When an agent states what a move will make possible, we can ask four questions in order. Is the statement recorded before the action? Is it accurate against the game state and the rules? Does the action agree with it? Does the behavior change when the size of the consequence changes? Work on LLM negotiation separates knowing a counterpart's position from using that knowledge~\cite{cosentino2026counterparty}, and work on planning tests reasoning about how actions change the world~\cite{valmeekam2023planbench}. We ask the same kind of question about two consequences in one environment: what a move opens for the agent itself, and what it opens for a rival.

The evidence is twenty complete four-player games of MonopolyBench played by one roster of four configured systems: GPT-6.1 Sol, Gemini 4 Argon, Claude Opus 5.5 and Grok 4.7. The records hold \fact{p.totals.decisions} decisions, \fact{o.census.responses.n} responses to trade offers and \fact{d.episodes} debts that required a liquidation menu. Each offer is tied to the state its recipient saw, and each debt to its first menu, the actions taken and the outcome.

\paragraph{The agent's own options.}
At the first menu of a debt, agents quit in \fact{d.above-bound.declared-at-first-menu} of the \fact{d.above-bound.episodes} debts that exceeded everything they owned and in \fact{x.debt.payable.declared} of the \fact{d.payable} debts that could be paid. Sixteen of the payable debts could be met only with mortgages that the menu would not list until buildings were sold. Agents paid fourteen. In thirteen of the fourteen, the text stored at the first decision counts those unlisted mortgages before any building is sold, and in the other it raises them as an open question. In both failures the text adds up the listed items only, and the agent declares bankruptcy. The prompt states neither the rule that a sale releases mortgages nor any mortgage value. The figures the agents state match the records in \fact{x.nc.numbers.equal} of \fact{x.nc.numbers.checked} checks.

\paragraph{A rival's options.}
Recipients accepted \fact{x.po.accept} of \fact{x.po.n} opening offers that would complete only the proposer's color group. They rejected \fact{x.swpo.rej.pct}\% of such swaps, against \fact{x.swn.rej.pct}\% of swaps that completed nobody's group. The gap has the same direction for every recipient, proposer and seed block, and its size depends on which proposers are counted. Cash bids for a completing deed were rejected less often (\fact{x.cbpo.rej.pct}\%) and countered about twice as often. Refusal texts mention the completed set in nearly every case. Gemini states a reciprocity rule, Grok states never-sell commitments with estimates of the rival's building cash, and Sol and Opus mostly name a price. Rejection does not rise with the number of houses the proposer could then afford.

\paragraph{Third parties and the uses of talk.}
Eight sourced incidents show these statements in context. They include a refusal of one's own monopoly to deny a rival's, a deed sold as insurance against a third player's inheritance, a ten-action payment that ends on the announced \$14, a bluff, a corrected false claim, a public call to withhold cash and a withdrawn assurance.

The two halves reach different levels of evidence. For the agent's own financing, the statements are prospective, mostly accurate and followed by matching actions, in cases fixed by rule arithmetic. For rivals, stated concern and matching refusals attach to the event of completion, with no measured response to the size of the rival's opportunity. Neither result measures whether the behavior was effective. All of it describes four configured systems in one cohort, two of which ran under altered response conditions (Section~\ref{sec:setting}).

The paper contributes a census of all \fact{d.episodes} debt episodes partitioned by what the first menu could raise, a complete table of the \fact{o.census.acceptance-legal} acceptance-legal offers by form, completion and stage with a stated sensitivity analysis, and eight incidents with exact quotations, channels, receipt checks and continuations.
